\chapter{Nhanh và Chậm}
\markboth{Nhanh và Chậm}{Fast and Slow}

\section{Một người muốn giàu nhanh nên mất cả vốn lẫn bình tĩnh.}
\begin{truyen}{Một người muốn giàu nhanh nên mất cả vốn lẫn bình tĩnh.}{Fast gains can create long losses.}
\chuhoa{V}{ì nôn nóng, anh lao vào những lời mời hấp dẫn và bỏ qua những câu hỏi căn bản. Mất tiền xong, anh mới hiểu nhiều con đường nhanh không phải ngắn hơn; chúng chỉ gãy sớm hơn.}
\textit{(Out of impatience, he rushed into attractive offers and ignored basic questions. After losing money, he understood that many fast roads are not shorter; they simply break sooner.)}
\end{truyen}
\begin{baihoc}
Đi chậm mà chắc thường không hấp dẫn bằng đi nhanh, nhưng ít để lại những hố sâu hơn.
\textit{(Slow and steady is less exciting than speed, but it leaves fewer deep pits behind.)}
\end{baihoc}
\begin{ghinhoanh}
\textbf{Short English to remember:}

Fast gains can create long losses.
\end{ghinhoanh}
\begin{tuhocnhanh}
1. Em đang vội ở chỗ nào chỉ vì muốn có kết quả sớm? \textit{(Where are you rushing simply to see results sooner?)}

2. Nếu đi chậm hơn một chút, em giữ được gì? \textit{(If you slowed down a little, what could you protect?)}
\end{tuhocnhanh}
\ngancach

\section{Một học sinh tiến bộ chậm nhưng rất chắc nền.}
\begin{truyen}{Một học sinh tiến bộ chậm nhưng rất chắc nền.}{Slow growth can be strong growth.}
\chuhoa{E}{m không nổi bật trong vài tháng đầu, nhưng học rất kỹ phần căn bản. Một năm sau, khi người khác bắt đầu hụt hơi, em lại tiến đều. Có những sự chậm chỉ là cách nền móng được xây sâu hơn.}
\textit{(He was not impressive in the first months, but studied the basics thoroughly. A year later, when others began to fade, he kept moving steadily. Some slowness is simply the way a foundation is built deeper.)}
\end{truyen}
\begin{baihoc}
Chậm không luôn là kém; nhiều khi nó là dấu hiệu của sự bền.
\textit{(Slow does not always mean weak; often it signals durability.)}
\end{baihoc}
\begin{ghinhoanh}
\textbf{Short English to remember:}

Slow growth can be strong growth.
\end{ghinhoanh}
\begin{tuhocnhanh}
1. Em có đang tự ép mình phải nhanh như người khác không? \textit{(Are you forcing yourself to move as fast as others?)}

2. Nền móng nào của em vẫn cần thêm thời gian? \textit{(Which foundation in you still needs more time?)}
\end{tuhocnhanh}
\ngancach

\section{Một người yêu quá nhanh rồi rời nhau cũng rất nhanh.}
\begin{truyen}{Một người yêu quá nhanh rồi rời nhau cũng rất nhanh.}{What starts too fast may not root deeply.}
\chuhoa{C}{ả hai cuốn vào nhau bằng cảm xúc mạnh, gần như không kịp hiểu nhau. Khi những khác biệt thật hiện ra, mối quan hệ rã nhanh như lúc nó bắt đầu. Có những điều đi chậm lại mới đủ chỗ cho hiểu biết lớn lên.}
\textit{(They were pulled together by strong emotions before truly knowing each other. When real differences appeared, the relationship ended as quickly as it began. Some things need slowness for understanding to grow.)}
\end{truyen}
\begin{baihoc}
Nhanh cho cảm xúc không đồng nghĩa với bền cho tình cảm.
\textit{(Fast emotion does not guarantee lasting affection.)}
\end{baihoc}
\begin{ghinhoanh}
\textbf{Short English to remember:}

Fast feelings are not deep roots.
\end{ghinhoanh}
\begin{tuhocnhanh}
1. Em có đang muốn điều gì lớn lên quá nhanh không? \textit{(Are you wanting something to grow too quickly?)}

2. Điều gì cần thêm thời gian để thành thật hơn? \textit{(What needs more time to become more real?)}
\end{tuhocnhanh}
\ngancach

\section{Một người luôn quyết rất nhanh rồi sửa rất lâu.}
\begin{truyen}{Một người luôn quyết rất nhanh rồi sửa rất lâu.}{Quick decisions can demand slow repair.}
\chuhoa{A}{nh tự hào mình quyết đoán, nhưng nhiều quyết định của anh chỉ là thiếu kiên nhẫn. Mỗi lần xử lý hậu quả, anh mới hiểu nhanh lúc chọn không giúp tiết kiệm thời gian nếu sau đó phải sửa quá lâu.}
\textit{(He was proud of being decisive, but many decisions were simply impatience. Each time he had to fix the consequences, he learned that quick choices do not save time if repair takes much longer.)}
\end{truyen}
\begin{baihoc}
Nhanh ở khâu chọn mà chậm ở khâu sửa chưa chắc đã là khôn.
\textit{(Being fast in choosing but slow in repairing is not necessarily wise.)}
\end{baihoc}
\begin{ghinhoanh}
\textbf{Short English to remember:}

Quick choices can demand slow repair.
\end{ghinhoanh}
\begin{tuhocnhanh}
1. Em có đang quyết nhanh ở chỗ cần suy nghĩ thêm không? \textit{(Are you deciding too quickly where more thought is needed?)}

2. Hậu quả nào em đang phải sửa vì một phút nóng vội? \textit{(What consequence are you repairing because of a rushed moment?)}
\end{tuhocnhanh}
\ngancach

\section{Người đi chậm nhưng biết dừng để nhìn lại mình.}
\begin{truyen}{Người đi chậm nhưng biết dừng để nhìn lại mình.}{Slow travel sees more truth.}
\chuhoa{K}{hông phải ai đi chậm cũng tụt lại. Có người chậm vì họ còn đủ tỉnh để tự hỏi mình đang đi đúng chưa, đang sống tử tế chưa. Chính những lần dừng ấy giữ cho đường đi không lạc quá xa.}
\textit{(Not everyone who moves slowly is falling behind. Some move slowly because they are awake enough to ask whether they are on the right road and living decently. Those pauses keep the journey from drifting too far.)}
\end{truyen}
\begin{baihoc}
Dừng lại để nhìn mình không phải chậm trễ; đó là một phần của trưởng thành.
\textit{(Stopping to examine yourself is not delay; it is part of maturity.)}
\end{baihoc}
\begin{ghinhoanh}
\textbf{Short English to remember:}

Slow travel sees more truth.
\end{ghinhoanh}
\begin{tuhocnhanh}
1. Lâu rồi em đã dừng lại để nhìn mình chưa? \textit{(How long has it been since you paused to look at yourself?)}

2. Nếu chậm lại hôm nay, em sẽ thấy điều gì rõ hơn? \textit{(If you slowed down today, what might become clearer?)}
\end{tuhocnhanh}
\ngancach